

%%%%%%%%%%%%%%%%%%%%%%
%%%%%%%%%%%%%%%%%%%%%%
\section{Methods}
%%%%%%%%%%%%%%%%%%%%%%
%%%%%%%%%%%%%%%%%%%%%%

Let $\mathcal{S}=\{\mathbf{s}_{1},\mathbf{s}_{2},\mathbf{s}_{3},...\}$ be a universe of protein sequences, where each $\mathbf{s}\in\mathcal{S}$ is a string of symbols from an alphabet of amino acids $\mathcal{A}=\left\{ A,C,D,\ldots,Y\right\} $. Let also $\mathcal{S}_{L}\subset\mathcal{S}$ be a set of labeled sequences, e.g. those with known structural class or function, that is provided as training data. The objective is to use supervised framework to probabilistically annotate the remaining sequences, i.e. those from the unlabeled set $\mathcal{S}_{U}=\mathcal{S}-\mathcal{S}_{L}$. 

To map protein sequences into a real-valued vector representation, let $\mathbf{s}=s_{1}s_{2}\ldots s_{\ell}$ be a length-$\ell$ protein sequence in $\mathcal{S}$ and $\mathbf{p}=(p_{1},p_{2},\ldots,p_{\ell})$ some property vector defined by any particular mapping from $\mathcal{A}^{\ell}$ to $\mathbb{R}^{\ell}$. For example, $\mathbf{p}$ may be provided as a vector of hydrophobicity indices corresponding to amino acids in $\mathbf{s}$. Alternatively, it can be represented as predicted helical propensities as outputted by some predictor of secondary structure. For those sequences with available structures, $\mathbf{p}$ may correspond to a sequence of dihedral angles calculated from the protein structure model of $\mathbf{s}$. 

Consider now a single property vector $\mathbf{p}$, such as hydrophobicity profile, corresponding to a particular sequence $\mathbf{s}\in\mathcal{S}$. We split $\mathbf{p}$ into $n$-dimensional sub-vectors $\mathbf{p}_{[1,n]}$, $\mathbf{p}_{[2,n+1]}$, $\ldots$, $\mathbf{p}_{[\ell-n+1,\ell]}$, where $\mathbf{p}_{[i,i+j]}=(p_{i},p_{i+1},...,p_{i+j})$, and $n\ll\ell$ is a small integer. For example, $\mathbf{p}_{[1,n]}$ corresponds to the first $n$ elements of $\mathbf{p}$ as shown in Figure \ref{fig:VQ_schematic}(a). For a property vector (amino acid sequence) of length $\ell$, there are $\ell-n+1$ length-$n$ sub-vectors.

Given a set of length-$n$ property sub-vectors $\mathcal{P}$ derived from the sequence universe $\mathcal{S}$, we generate a partition of $\mathbb{R}^{n}$ into $m$ regions $\mathcal{R}=\left\{ R_{1},R_{2},\ldots,R_{m}\right\} $. These regions are represented by a set of $n$-dimensional vectors, or centroids, $\mathcal{C}=\left\{ \mathbf{c}_{1},\mathbf{c}_{2},\ldots,\mathbf{c}_{m}\right\} $. Each region $R_{i}$ represents a Voronoi region such that 

\[
R_{i}=\left\{ \mathbf{x}:d(\mathbf{x},\mathbf{c}_{i})\leq d(\mathbf{x},\mathbf{c}_{j}), j \neq i \right\} ,
\]

\noindent where $d(\mathbf{x},\mathbf{c})$ is the Euclidean distance between vector $\mathbf{x}$ and centroid $\mathbf{c}$. We determine $\mathcal{C}$ using k-means clustering, where the initial set of clusters is generated by the splitting method \citep{gersho1992vector}.


%%%%%%%%%%%%%%%%%%%%%%
%%%%%%%%%%%%%%%%%%%%%%
\subsection{Count-based property kernels}
\label{sec:propertykernels}
%%%%%%%%%%%%%%%%%%%%%%
%%%%%%%%%%%%%%%%%%%%%%

Variable length property vectors can be transformed into fixed $m$-length vectors using the partition of $\mathbb{R}^{n}$ defined by $\mathcal{R}$. Specifically, a property vector $\mathbf{p}$ is mapped into a vector of length $m$ as 

\[
\mathbf{x}=\left(\varphi_{1}(\mathbf{p}),\varphi_{2}(\mathbf{p}),\ldots,\varphi_{m}(\mathbf{p})\right),
\]

\noindent where $\varphi_{i}(\mathbf{p})$ is the number of $n$-dimensional vectors $\mathbf{p}_{[.]}$ in $\mathbf{p}$ that belong to partition $R_{i}$. Given two property vectors $\mathbf{p}$ and $\mathbf{q}$ and their respective count vectors $\mathbf{x}$ and $\mathbf{y}$, a \textit{count-based property kernel} is defined as

\[
k(\mathbf{p},\mathbf{q})=\mathbf{x}^{T}\mathbf{y},
\]
where $T$ is the transpose operator. Note that in this notation each count vector is assumed to be a column vector, i.e. $\left(a,b,c\right)=\left[a\; b\; c\right]^{T}$, as in \cite{strang2003introduction}. Since the function $k(\mathbf{p},\mathbf{q})$ is defined as an inner product between two count vectors, it is a kernel function \citep{Haussler1999}.

%%%%%%%%%%%%%%%%%%%%%%
%%%%%%%%%%%%%%%%%%%%%%
%\subsection{Similarity augmented count based property kernel}
%\label{sec:distancekernel}
%%%%%%%%%%%%%%%%%%%%%%
%%%%%%%%%%%%%%%%%%%%%%

The count-based property kernels can be generalized by utilizing similarities between centroids to provide partial counts to each of the regions in $\mathcal{R}$. First, a similarity function between any two centroids is defined as 

\[
s(\mathbf{c}_{i},\mathbf{c}_{j})=e^{-\alpha\cdot\frac{d(\mathbf{c}_{i},\mathbf{c}_{j})}{d_{\max}}},
\]

\noindent where $d(\mathbf{c}_{i},\mathbf{c}_{j})$ is the Euclidean distance between $\mathbf{c}_{i}$ and $\mathbf{c}_{j}$ and $d_{\max}$ is the maximum distance between any two centroids in $\mathcal{C}$. The parameter $\alpha$ is a positive constant that defines an extent to which large distances are penalized by the non-linear mapping above, and is optimized in the model evaluation step. A generalized property kernel can now be defined as

\[
k_{g}(\mathbf{p},\mathbf{q})=(\mathbf{Sx})^{T}(\mathbf{Sy})=\mathbf{x}^{T}\mathbf{S}^{T}\mathbf{S}\mathbf{y}=\mathbf{x}^{T}\mathbf{Q}\mathbf{y},
\]

\noindent where elements of the similarity matrix $\mathbf{S}$ are defined as $s_{ij}=s(\mathbf{c}_{i},\mathbf{c}_{j})$. Clearly, for $\mathbf{S}=\mathbf{I}$ the generalized kernel is equivalent to the basic count-based property kernel. If \textbf{$\mathbf{S}$} is a non-singular matrix, it is easy to show that $\mathbf{Q}$ is a positive definite matrix \citep{strang2003introduction}. Therefore, $k_{g}(\mathbf{p},\mathbf{q})$ is a kernel function. Note that the generalized count-based kernel does not lead to a sparse representation, even for large $m$. 


%%%%%%%%%%%%%%%%%%%%%%
%%%%%%%%%%%%%%%%%%%%%%
\subsection{Additive property kernels}
\label{sec:additivekernels}
%%%%%%%%%%%%%%%%%%%%%%
%%%%%%%%%%%%%%%%%%%%%%

Let us again consider vector $\mathbf{p}$ and partition $\mathcal{R}$. We map $\mathbf{p}$ into an $m$-dimensional as 

\[
\mathbf{x}=\left(\psi_{1}(\mathbf{p}),\psi_{2}(\mathbf{p}),\ldots,\psi_{m}(\mathbf{p})\right),
\]
where 

\[
\psi_{i}(\mathbf{p})=\sum_{j=1}^{\ell-n+1}s\left(\mathbf{p}_{\left[j,j+n-1\right]},\mathbf{c}_{i}\right),
\]
is the cumulative similarity between vectors $\mathbf{p}_{[.]}$ in $\mathbf{p}$ and centroid $\mathbf{c}_{i}$, and $s\left(\cdot\right)$ is the vector similarity function defined above. Given two property vectors $\mathbf{p}$ and $\mathbf{q}$ and their respective count vectors $\mathbf{x}$ and $\mathbf{y}$, we now define an \textit{additive property kernel} as $k(\mathbf{p},\mathbf{q})=\mathbf{x}^{T}\mathbf{y}$. 


%%%%%%%%%%%%%%%%%%%%%%
%%%%%%%%%%%%%%%%%%%%%%
\subsection{Combined kernel function}
\label{sec:combinedkernels}
%%%%%%%%%%%%%%%%%%%%%%
%%%%%%%%%%%%%%%%%%%%%%

Given a set of property kernels $\left\{ k_{i}(\mathbf{x},\mathbf{y})\right\} $, we construct the composite property kernel as a linear combination

\[
k(\mathbf{x}, \mathbf{y})=\sum_{i}k_{i}(\mathbf{x},\mathbf{y})
\]

%Let $k_{hy}(\mathbf{x},\mathbf{y})$, $k_{fl}(\mathbf{x},\mathbf{y})$, $k_{he}(\mathbf{x},\mathbf{y})$, $k_{sh}(\mathbf{x},\mathbf{y})$, $k_{id}(\mathbf{x},\mathbf{y})$, be the property kernels representing hydrophobicity, flexibility, helix propensity, sheet propensity, and intrinsic disorder, respectively. The composite property kernel is calculated as a linear combination

%\[
%k(\mathbf{x},\mathbf{y})=k_{hy}(\mathbf{x},\mathbf{y})+k_{fl}(\mathbf{x},\mathbf{y})+k_{he}(\mathbf{x},\mathbf{y})+k_{sh}(\mathbf{x},\mathbf{y})+k_{id}(\mathbf{x},\mathbf{y}),
%\]

\noindent where before and after combining, each kernel is normalized using

\[
k(\mathbf{x},\mathbf{y})\leftarrow\frac{k(\mathbf{x},\mathbf{y})}{\sqrt{k(\mathbf{x},\mathbf{x})k(\mathbf{y},\mathbf{y})}}.
\]


%%%%%%%%%%%%%%%%%%%%%%
%%%%%%%%%%%%%%%%%%%%%%
\subsection{Computational complexity}
%%%%%%%%%%%%%%%%%%%%%%
%%%%%%%%%%%%%%%%%%%%%%

The computation of the count vector can be accomplished in $O(\ell mn)$ time if each $n$-dimensional vector from $\mathbf{p}$ is compared with all centroids in $\mathcal{C}$. Approximation algorithms are available through a decision tree like organization of the centroids. In such a case, only $\log m$ operations are needed resulting in $O(\ell n\log m)$ time \citep{gersho1992vector}; however, there is no guarantee that the closest centroid will be found. Using a full-search scheme, the basic and generalized property kernels can be computed in $O(\ell nm)$ and $O(m^{2}+\ell nm)$ times, respectively. The storage requirements include $O(mn)$ for storing $\mathcal{C}$ and $O(m^{2})$ for storing \textbf{$\mathbf{Q}$}. The additive kernel has complexity $O(\ell nm)$.

%%%%%%%%%%%%%%%%%%%%%%
%%%%%%%%%%%%%%%%%%%%%%
\subsection{String kernel}
%%%%%%%%%%%%%%%%%%%%%%
%%%%%%%%%%%%%%%%%%%%%%

As a baseline method with which we could compare results from our VQ model we utilized a string kernel as described by \cite{Leslie2002} for a wide range of word sizes ($n\in \{1\dots 10\}$).  For a given word size $n$, a sparse $20^{n}$-length vector was created for a protein, where each value represented the number of times a potential $n$ length substring that could be generated using the 20 amino acid alphabet occurred.  An $\ell$-length sequence, $\mathbf{s}$, would contain $\ell-n+1$ such overlapping strings.



